\section{Cartan Comparison: Residuals and Normal Source Compatibility}
\label{v490:cartan-source}

This section records a new research result, not an identification of the
native MM/HC action. It follows the action and geometry distinctions in
Section~\ref{v480:hc-dynamics-geometry}. The model below is an explicitly
added neutral Einstein--Cartan comparison with three complex commuting
spinors. Its finite jets do not constitute a solution on an open set.
Neither its matrices nor its six source coordinates are identified with
matter, photon or tachyon frames, C7 carriers, physical masses or a
measured gravitational coupling.

\subsection{Declared Model and Residual Convention}
\label{v490:cartan-model}

Take independent fields: an invertible real coframe $E^a{}_{\mu}$, a real
Lorentz connection $\omega_{\mu,ab}=-\omega_{\mu,ba}$, the spinor triple
$\Psi$, and a real multiplier density $\lambda_a{}^\mu$ of coordinate
weight one. Put $v=|\det E|$, $e=E^{-1}$ and
$\eta=\operatorname{diag}(-1,1,1,1)$. In the declared Dirac convention,
$K^a=\beta\Gamma^a$ and
\[
 B_\mu=-\tfrac14\omega_{\mu,ab}\Gamma^a\Gamma^b,
 \qquad D_\mu\Psi=\partial_\mu\Psi+B_\mu\Psi,
 \qquad C^a{}_{\mu}=\operatorname{Im}
       (\Psi^\dagger QK^aD_\mu\Psi).
\]
Here $Q$ acts on the three channels and $K^a$ on spin indices; tensor
products are implicit. A dagger is the ordinary complex adjoint, not a
new native comparison map. The assumed action is
\begin{align}
 S_{\rm cmp}&=\int d^4x\,
 \left\{v\left(\frac c2 R(E,\omega)-u\right)
             +\mathcal L_D+\lambda_a{}^\mu(E^a{}_{\mu}-C^a{}_{\mu})\right\},
 \label{v490:cartan-action}\\
 \mathcal L_D&=-v\operatorname{Im}
  (\Psi^\dagger e_a{}^\mu K^aD_\mu\Psi)
  -v\operatorname{Re}(\Psi^\dagger(M\otimes\beta)\Psi).
 \nonumber
\end{align}
The numerical branch uses $c=1$, $u=0$ and fixed dimensionless comparison
matrices
\[
 M=\begin{pmatrix}0.6&0&4\\0&0.6&0.7\\4&0.7&0.6\end{pmatrix},
 \qquad Q=\frac i2(E_{12}-E_{21}),
\]
where $E_{jk}$ denotes a channel matrix unit. These choices are not a
prediction of couplings or a derivation of the Einstein term. No internal
charged equation is suppressed: neutrality is part of this comparison.

Write $F=C-E$, and denote the independent coframe and full real-action
probe Euler densities by $\mathfrak E_a{}^\mu$ and $\mathcal R$. Thus the
reduced variation after exact algebraic Cartan elimination is
\[
 \delta S_{\rm cmp}=\int d^4x\,
 \left\{\mathfrak E_a{}^\mu\delta E^a{}_{\mu}
       +2\operatorname{Re}(\delta\Psi^\dagger\mathcal R)
       -F^a{}_{\mu}\delta\lambda_a{}^\mu\right\}.
\]
The connection solution must be differentiated compatibly. Setting only
its point Euler value to zero while selecting unrelated connection
derivatives does not define this reduced field theory. Its sourced
Cartan curvature is not generally the Levi--Civita curvature of $E$.

\subsection{Ward Identities and a Conditional Subsidiary System}
\label{v490:cartan-ward}

For compactly supported coordinate variations, the spinor is a coordinate
scalar and $\lambda$ is a vector density. Its variation includes
$\lambda_a{}^\mu\partial_\nu\xi^\nu$, in addition to transport and
contravariant-index terms. Integration by parts therefore gives
\begin{align}
0={}&E^a{}_{\nu}\partial_\mu\mathfrak E_a{}^\mu
 +\mathfrak E_a{}^\mu(\partial_\mu E^a{}_{\nu}
                             -\partial_\nu E^a{}_{\mu})\nonumber\\
 &+\lambda_a{}^\mu(\partial_\mu F^a{}_{\nu}
                             -\partial_\nu F^a{}_{\mu})
 +F^a{}_{\nu}\partial_\mu\lambda_a{}^\mu
 -2\operatorname{Re}((\partial_\nu\Psi)^\dagger\mathcal R).
 \label{v490:coordinate-identity}
\end{align}
For $J^T\eta+\eta J=0$ and
$S_J=-(\eta J)_{ab}\Gamma^a\Gamma^b/4$, Lorentz invariance similarly gives
\[
 \mathfrak E:(JE)+2\operatorname{Re}((S_J\Psi)^\dagger\mathcal R)
                          +F:(J^T\lambda)=0.
\]
These identities retain all residual sources. Their source-free
specialization requires $F=\mathcal R=0$ as field equations on an open
region, not merely at the prepared point.

\begin{proposition}[Conditional gravitational constraint propagation]
\label{v490:subsidiary}
Assume exact compatible Cartan elimination, $F=\mathcal R=0$ on an open
region, a sufficiently smooth Lorentzian metric with spacelike slices,
and the spatially projected trace-reversed coframe Euler evolution
equations. With zero initial gravitational constraints and no incoming
constraint flux, the four gravitational constraints remain zero.
This statement does not propagate $F$ or $\mathcal R$ from initial data.
\end{proposition}
\begin{proof}
The two identities make the coframe Euler tensor density symmetric and
Levi--Civita divergence-free under the stated hypotheses. Divide out the
volume density to obtain a symmetric residual tensor $T^{\mu\nu}$ with
$\nabla_\mu T^{\mu\nu}=0$. For the unit slice normal $n$, put
$H=T(n,n)$ and $M_i=-T(n,\partial_i)$. The spatial projection of
$T^{\mu\nu}-g^{\mu\nu}\operatorname{tr}_g T/2$ vanishes by the chosen
evolution equations and gives spatial stress $P^{ij}=Hh^{ij}$.
In lapse $N>0$, shift $\beta^i$ and positive slice metric $h_{ij}$, the
principal symbol for $U=(H,M_i)$ is
\[
 A(\xi)=-(\beta\cdot\xi)I_4+
 N\begin{pmatrix}0&(h^{-1}\xi)^T\\\xi&0_{3\times3}\end{pmatrix},
 \qquad W=\operatorname{diag}(1,h^{-1})>0.
\]
Since $WA(\xi)$ is symmetric, the energy
$\tfrac12\int(H^2+h^{ij}M_iM_j)\sqrt{\det h}\,d^3x$ obeys
$\mathcal E'(t)\leq C(t)\mathcal E(t)$ with bounded coefficient
derivatives and no incoming flux. Gronwall's inequality proves the claim.
With nonzero shift the evolution condition is the geometric spatial
projection, not simply the vanishing of contravariant coordinate $ij$
components. The full lower-order lapse, shift and connection terms must
be retained in the energy estimate.
\end{proof}

The saved Taylor fields do not satisfy the open-set hypotheses of this
proposition. In particular, the constitutive identity
\[
 \partial_0F^a{}_{i}=\partial_iF^a{}_{0}
              +(\partial_0F^a{}_{i}-\partial_iF^a{}_{0})
\]
does not close evolution unless its mixed-curl source is controlled.
Ward identities alone are not a full coupled propagation theorem.

\subsection{Finite Normal Reduction with Exposed Source Coordinates}
\label{v490:normal-reduction}

At the prepared normal point $E=I$, $\partial E=0$, hold the lower and
spatial data and all action coefficients fixed. The real normal unknowns
are
\[
 x=(E_{,00}\ [16],\ \operatorname{realify}\Psi_{,0}\ [24],
                    \lambda_{,0}\ [16])\in\mathbb R^{56}.
\]
The 16 constitutive, 24 real probe and 16 coframe Euler rows, together
with twelve mixed constitutive curls, form an affine residual
$b+Nx\in\mathbb R^{68}$. The curls are necessary consequences of
$F=0$ as a field equation, not twelve independent additional action laws.
The numerical rank rises from 34 for the original rows to 40 with the
curls. Its kernel has an exhibited ten-dimensional leading normal gauge
span and a six-dimensional complementary span. This is not a complete
regional gauge reduction.

\begin{proposition}[Compatible finite normal inverse]
\label{v490:normal-inverse}
Assume $\operatorname{rank}N=40$, and orthonormal matrices
$G\in\mathbb R^{56\times10}$, $U\in\mathbb R^{56\times6}$ with
$\ker N=\operatorname{im}G\mathbin{\oplus}\operatorname{im}U$.
Let $L_N$ span $\ker N^T$ orthonormally. Then $Nx=r$ is solvable if and
only if $L_N^Tr=0$. Every solution is uniquely parametrized by
\[
 x=N^+r+Gg+Uq,\qquad g\in\mathbb R^{10},\quad q\in\mathbb R^6.
\]
\end{proposition}
\begin{proof}
$\operatorname{im}N=(\ker N^T)^\perp$ gives compatibility.
For compatible $r$, $NN^+r=r$; subtracting this particular solution
leaves exactly $\ker N$. Orthogonality gives unique $g,q$.
For incompatible $r$ the same expression leaves residual $-L_NL_N^Tr$;
a least-squares output is not an exact solution.
\end{proof}

There are 28 left compatibility rows on arbitrary $r$, not necessarily
28 independent physical field laws. For increments about the saved jet,
$r$ is the target residual minus its saved value, rather than the absolute
right-hand side $-b$. The chosen $Uq$ changes only $E_{,00}$ and the
spatial-column entries $\lambda_{a,0}{}^i$, leaving all $\Psi_{,0}$
and $\lambda_{a,0}{}^0$ fixed. These six coordinates are neither the
earlier mixed-curl defects nor established gauge or physical modes.

\subsection{Determination from Compatible Tangential Jets}
\label{v490:source-law}

Use the same current saved joint jet throughout. Let $q$ now be a source
increment with gauge and normal residual target fixed. Let $z$ contain
increments of $(\Psi_{,2},\lambda_{,2},E_{,3},\Psi_{,3})$ in the real
unique-symmetric convention, of dimensions $(240,160,320,480)$.
Direct differentiation of the current equations gives the jointly affine
384-row residual
\begin{equation}
 \mathcal H(q,z)=h_0+Bq+Az.
 \label{v490:higher-affine}
\end{equation}
Its blocks are first constitutive derivatives (64), first real probe
derivatives (96), first coframe Euler derivatives (64), and second
constitutive derivatives in unique pairs (160). Affinity is specific to
this source variation and order: the lower spinor and its first jet are
fixed, so variable factors multiply fixed lower factors. It does not
extend automatically to later derivative equations.

For all 1200 higher coordinates allowed, $\operatorname{rank}A=332$ and
$\operatorname{rank}(L_A^TB)=0$ numerically, where $L_A$ spans
$\ker A^T$. Thus the frozen-higher rank-six response of $B$ can be
absorbed by higher increments at this order. It does not fix $q$.

Instead preserve the represented finite tangential germ of $E$, $E_{,0}$,
$\Psi$ and $\lambda$ on $x^0=0$. Split $z$ into free coordinates $w$ and
prescribed tangential coordinates $d$. With $i,j,k\in\{1,2,3\}$, the
free slots and real dimensions are
\[
 \begin{array}{c|c|r}
 \text{block}&\text{allowed derivative indices}&\text{dimension}\\\hline
 \Psi_{,2}&\text{at least one }0&96\\
 \lambda_{,2}&\text{at least one }0&64\\
 E_{,3}&\text{at least two }0&64\\
 \Psi_{,3}&\text{at least one }0&240
 \end{array}
\]
Thus $w\in\mathbb R^{464}$ and $d\in\mathbb R^{736}$. In particular
$E_{,0ij}$ is prescribed because it differentiates the fixed $E_{,0}$
trace tangentially. This is a finite germ at a point, not a construction
of admissible surface functions. Write $Az=A_fw+A_dd$.

\begin{proposition}[Conditional source law for prescribed tangential data]
\label{v490:source-law-theorem}
Let $L_f$ span $\ker A_f^T$ orthonormally, set
\[
 J=L_f^TB,\qquad b(d)=-L_f^T(h_0+A_dd),
\]
and assume $\operatorname{rank}J=6$. A source and higher lift exist if
and only if $(I-JJ^+)b(d)=0$. For compatible $d$ the source is unique:
\[
 q(d)=J^+b(d),\qquad
 w=-A_f^+(h_0+A_dd+Bq(d))+k,\quad k\in\ker A_f.
\]
\end{proposition}
\begin{proof}
Eliminating $w$ requires and suffices for
$L_f^T(h_0+A_dd+Bq)=0$, namely $Jq=b(d)$.
Full column rank makes $q$ unique if this equation is consistent.
The orthogonal projector $JJ^+$ tests that consistency; the remaining
formula gives exactly all higher lifts.
\end{proof}

The numerical response for this mask has rank 200 and the six singular
values of $J$ are approximately
\[
 1.89152,\ 1.22410,\ 0.48769,\ 0.45652,\ 0.26352,\ 0.13424.
\]
The residual compatibility map on $d$ has numerical rank 126. This is a
finite matrix rank, not a count of physical constraints. All ranks are
pointwise floating calculations with declared cutoffs, not certified
regional bounds. Incompatible prescribed data are rejected, not silently
projected onto admissible ones. In particular $q=0$ is not an added law;
near-zero increments for unchanged data refer to the saved germ only.

\subsection{Reproducibility, Negative Control and Remaining Obligations}
\label{v490:cartan-evidence}

Four compatible-data controls use
$q=a(1,-0.5,0.3,0.2,0.4,-0.1)$ with
$a\in\{0,0.01,-0.01,0.02\}$. An unrestricted compatible lift supplies
the prescribed $d$, which is then held exactly fixed while the conditional
law recovers $q$. Recovery errors are below $2.5\times10^{-15}$;
direct residual norms are below $2.4\times10^{-14}$ for all 68 normal
rows and $3.7\times10^{-11}$ for all 384 higher rows. These controls
verify a known compatible family, not arbitrary surface admissibility.
The computed $\|h_0\|\simeq4.78\times10^{-11}$ is retained in the law.

A separate saved completion satisfies all 1114 packed rows of the later
joint hierarchy to residual about $5.38\times10^{-10}$. The new source
lifts do not automatically preserve it. For $a=0.01$, retaining its old
highest $E_{,4},\lambda_{,3},\Psi_{,4}$ gives a later residual about
98.96, with third constitutive block about 98.60, even after recomputing
the changed Levi--Civita cubic. This is a failed frozen-highest control,
not a nonexistence proof for a re-prepared completion. The earlier
coefficient-domain neighborhood estimate does not transfer to the new
witnesses without recalculation.

The source files and pinned numerical reports are
\nolinkurl{analysis/HC_Cartan_Normal_Source_Reduction.py} and
\nolinkurl{analysis/HC_Cartan_Normal_Source_Compatibility.py}, with
their matching JSON reports and derivation notes. The latter module uses
the fingerprint-checked current joint loader, not the older regular-point
completion. Sixteen source-compatibility and thirteen normal-reduction
tests supplement the full audit; a test count is not a proof certificate.

For a future smooth normal chart
$x=x_{\rm part}+Gg+Uq$, spatial differentiation must retain
\[
 \partial_i x=\partial_i x_{\rm part}+(\partial_iG)g+G\partial_i g
                              +(\partial_iU)q+U\partial_iq.
\]
A QR/SVD frame at one point supplies neither these derivatives nor
arbitrary source functions. The next obligations are compatible surface
data, higher re-preparation with the prescribed tangential germ retained,
a smooth gauge/source chart, and a closed full evolution and subsidiary
estimate. No regional existence, full hyperbolicity, native-action
identification, OP5 closure or physical scale follows from the finite
source law alone.

\clearpage
